% Arduino入门教材

\documentclass[12pt,UTF8]{ctexbook}

% 设置纸张信息。
\input{../common/page}

% 设置字体，并解决显示难检字问题。
\xeCJKsetup{AutoFallBack=true}
% 注意：ctexbook类已默认设置SimSun为CJKrmdefault
% 以下设置用于确保字体回退和扩展字体可用
\setCJKfamilyfont{hei}{SimHei}
\setCJKfamilyfont{kai}{KaiTi}

% 目录 chapter 级别加点（.）。
\usepackage{titletoc}
\titlecontents{chapter}[0pt]{\vspace{3mm}\bf\addvspace{2pt}\filright}{\contentspush{\thecontentslabel\hspace{0.8em}}}{}{\titlerule*[8pt]{.}\contentspage}

% 支持目录点击跳转
\usepackage[colorlinks,linkcolor=blue,citecolor=blue,urlcolor=blue]{hyperref}

% 设置 part 和 chapter 标题格式。
\ctexset{
	part/name= {第,篇},
	part/number={\chinese{part}},
	chapter/name={第,章},
	chapter/number={\arabic{chapter}}
}

% 图片相关设置。
\usepackage{graphicx}
\graphicspath{{Images/}}

% 代码列表支持
\usepackage{listings}
\usepackage{xcolor}

\lstset{
	basicstyle=\ttfamily\small,
	keywordstyle=\color{blue},
	commentstyle=\color{green!50!black},
	stringstyle=\color{red},
	numbers=left,
	numberstyle=\tiny\color{gray},
	numbersep=5pt,
	frame=single,
	tabsize=4,
	breaklines=true,
	breakatwhitespace=false,
	showspaces=false,
	showstringspaces=false,
	showtabs=false,
	extendedchars=false,
	columns=fixed,
	language=C++,
	morekeywords={setup,loop,pinMode,digitalWrite,digitalRead,analogWrite,analogRead,delay,Serial,begin,println,print,available,read,attachInterrupt,detachInterrupt,noInterrupts,interrupts,millis,micros}
}

% 设置段落间距
\setlength{\parskip}{0.8em plus 0.2em minus 0.1em}

% 列表格式支持，列表项向右偏移。
\usepackage{enumitem}

\title{\heiti\zihao{0} Arduino}
\author{WangFei}
\date{\today}

\begin{document}

\maketitle
\tableofcontents

\frontmatter
\chapter{前言}

\mainmatter

%========== 第一篇：基础概念 ==========
\part{基础概念}

\chapter{Arduino简介}

\section{什么是Arduino}

Arduino 是一个开源的电子原型平台，它由两部分组成：

\begin{enumerate}
    \item \textbf{硬件}：一块带有微控制器的电路板（开发板）
    \item \textbf{软件}：一个用于编写和上传代码的集成开发环境（IDE）
\end{enumerate}

简单来说，Arduino 就是一个让您可以轻松控制电子设备的工具。您可以把它想象成一个"电子大脑"，通过编写程序来控制各种电子元件，如 LED 灯、电机、传感器等。

\section{Arduino的诞生与发展}

Arduino 项目始于 2005 年，由意大利米兰交互设计学院的 Massimo Banzi 和 David Cuartielles 等人创建。最初是为了帮助设计专业的学生能够快速原型制作而开发的。

由于其简单易用、价格低廉且开源的特点，Arduino 迅速在全球范围内流行起来。现在，它不仅被用于教育和艺术创作，还广泛应用于工业自动化、物联网、机器人等领域。

\section{Arduino能做什么}

Arduino 的应用范围非常广泛，以下是一些常见的应用场景：

\begin{itemize}
    \item \textbf{互动艺术}：制作灯光装置、声音艺术作品
    \item \textbf{智能家居}：控制灯光、温度、窗帘等
    \item \textbf{机器人}：制作遥控小车、机械臂等
    \item \textbf{传感器网络}：环境监测、数据采集
    \item \textbf{原型验证}：快速验证电子产品想法
\end{itemize}

只要您能想到的电子创意，Arduino 都能帮您实现！

\section{Arduino生态系统}

Arduino 不仅仅是一块开发板，它是一个完整的生态系统：

\begin{enumerate}
    \item \textbf{开发板}：多种型号的 Arduino 开发板，如 UNO、Nano、Mega 等
    \item \textbf{扩展板（Shields）}：可以堆叠在开发板上的功能扩展模块
    \item \textbf{传感器和模块}：各种传感器（温度、湿度、光线等）和执行器（电机、继电器等）
    \item \textbf{软件库}：丰富的开源代码库，方便实现各种功能
    \item \textbf{社区}：庞大的用户社区，有大量教程和项目分享
\end{enumerate}

\section{为什么选择Arduino}

相比其他电子开发平台，Arduino 有以下优势：

\begin{itemize}
    \item \textbf{简单易学}：编程语法简单，非常适合初学者
    \item \textbf{价格实惠}：开发板价格相对较低
    \item \textbf{开源}：硬件和软件都是开源的，可以自由修改
    \item \textbf{跨平台}：支持 Windows、Mac、Linux 系统
    \item \textbf{扩展性强}：可以连接各种传感器和模块
\end{itemize}

\chapter{Arduino硬件平台}

\section{Arduino开发板概述}

Arduino 开发板是整个 Arduino 系统的核心。它主要包含以下几个部分：

\begin{enumerate}
    \item \textbf{微控制器}：开发板的"大脑"，负责执行程序
    \item \textbf{电源接口}：为开发板供电
    \item \textbf{数字引脚}：用于数字信号的输入和输出
    \item \textbf{模拟引脚}：用于模拟信号的输入
    \item \textbf{USB接口}：用于连接电脑，上传程序和通信
    \item \textbf{LED指示灯}：电源指示灯和内置LED
\end{enumerate}

\section{常见Arduino开发板介绍}

市面上有多种 Arduino 开发板，每种都有其特点：

\begin{enumerate}
    \item \textbf{Arduino UNO}：最经典的型号，适合初学者
    \item \textbf{Arduino Nano}：体积小巧，适合空间有限的项目
    \item \textbf{Arduino Mega}：拥有更多引脚和更大内存，适合复杂项目
    \item \textbf{Arduino Leonardo}：支持 USB HID，可以模拟键盘和鼠标
    \item \textbf{Arduino Zero}：基于 ARM 处理器，性能更强
\end{enumerate}

\section{Arduino UNO详解}

Arduino UNO 是最适合初学者的开发板，让我们详细了解一下它的结构：

\begin{enumerate}
    \item \textbf{电源部分}：
    \begin{itemize}
        \item USB 接口：通过 USB 线连接电脑供电或上传程序
        \item DC 电源插孔：使用 9V 电池或电源适配器供电
        \item 电源开关和指示灯
    \end{itemize}
    
    \item \textbf{微控制器}：
    \begin{itemize}
        \item 型号：ATmega328P
        \item 工作电压：5V
        \item 数字 I/O 引脚：14 个（其中 6 个支持 PWM）
        \item 模拟输入引脚：6 个
        \item 内存：32KB Flash（用于存储程序），2KB SRAM，1KB EEPROM
    \end{itemize}
    
    \item \textbf{引脚说明}：
    \begin{itemize}
        \item 数字引脚（0-13）：用于数字信号输入输出
        \item 模拟引脚（A0-A5）：用于读取模拟信号
        \item 引脚旁边带有波浪线（~）的支持 PWM 输出
    \end{itemize}
\end{enumerate}

\section{其他型号简介}

\begin{enumerate}
    \item \textbf{Arduino Nano}：
    \begin{itemize}
        \item 体积小，适合面包板项目
        \item 与 UNO 使用相同的 ATmega328P 芯片
        \item 通过 mini-USB 接口连接
    \end{itemize}
    
    \item \textbf{Arduino Mega}：
    \begin{itemize}
        \item 更多引脚：54 个数字引脚，16 个模拟引脚
        \item 使用 ATmega2560 芯片
        \item 适合需要大量 I/O 的复杂项目
    \end{itemize}
    
    \item \textbf{Arduino Leonardo}：
    \begin{itemize}
        \item 使用 ATmega32U4 芯片
        \item 支持 USB HID 协议
        \item 可以模拟键盘、鼠标等输入设备
    \end{itemize}
\end{enumerate}

\section{Arduino扩展板介绍}

扩展板（Shields）是可以直接堆叠在 Arduino 开发板上的电路板，用于扩展功能：

\begin{enumerate}
    \item \textbf{电机扩展板}：用于控制直流电机和舵机
    \item \textbf{以太网扩展板}：让 Arduino 连接到互联网
    \item \textbf{蓝牙/WiFi 扩展板}：实现无线通信
    \item \textbf{LCD 显示屏扩展板}：显示文字和图形
    \item \textbf{传感器扩展板}：集成多种传感器
\end{enumerate}

\chapter{开发环境搭建}

\section{Arduino IDE安装}

Arduino IDE 是官方提供的集成开发环境，用于编写、编译和上传程序。

\begin{enumerate}
    \item \textbf{下载地址}：访问 Arduino 官方网站（https://www.arduino.cc/en/software）
    \item \textbf{安装步骤}：
    \begin{enumerate}
        \item 下载对应操作系统的安装包
        \item 运行安装程序，按照提示完成安装
        \item 启动 Arduino IDE
    \end{enumerate}
\end{enumerate}

\section{IDE界面介绍}

Arduino IDE 界面主要分为以下几个区域：

\begin{enumerate}
    \item \textbf{菜单栏}：包含文件、编辑、项目、工具、帮助等菜单
    \item \textbf{工具栏}：常用操作按钮（新建、打开、保存、上传等）
    \item \textbf{代码编辑区}：编写 Arduino 代码的地方
    \item \textbf{状态栏}：显示当前选中的开发板和端口信息
    \item \textbf{串口监视器}：用于查看串口输出信息
\end{enumerate}

\section{驱动安装}

某些 Arduino 开发板需要安装驱动程序才能被电脑识别：

\begin{enumerate}
    \item \textbf{连接开发板}：用 USB 线将 Arduino 连接到电脑
    \item \textbf{安装驱动}：
    \begin{itemize}
        \item Windows：系统会自动搜索驱动，可能需要手动安装
        \item Mac/Linux：通常无需额外驱动
    \end{itemize}
    \item \textbf{确认连接}：在 IDE 的"工具"-> "端口"中可以看到连接的串口
\end{enumerate}

\section{第一个程序：Blink}

让我们编写第一个 Arduino 程序 - 让 LED 闪烁：

\begin{lstlisting}
// 定义LED引脚
const int ledPin = 13;

void setup() {
  // 将LED引脚设置为输出模式
  pinMode(ledPin, OUTPUT);
}

void loop() {
  // 点亮LED
  digitalWrite(ledPin, HIGH);
  // 等待1秒
  delay(1000);
  // 熄灭LED
  digitalWrite(ledPin, LOW);
  // 等待1秒
  delay(1000);
}
\end{lstlisting}

上传步骤：
\begin{enumerate}
    \item 选择开发板："工具" -> "开发板" -> "Arduino UNO"
    \item 选择端口："工具" -> "端口" -> 选择正确的串口
    \item 点击"上传"按钮（右箭头图标）
\end{enumerate}

上传成功后，您会看到 Arduino 板上的内置 LED（第 13 引脚）开始闪烁。

\section{常见问题排查}

如果遇到问题，可以按照以下步骤排查：

\begin{enumerate}
    \item \textbf{检查连接}：确保 USB 线连接牢固
    \item \textbf{选择正确的开发板}：确认选择的开发板型号与实际一致
    \item \textbf{选择正确的端口}：在设备管理器中查看串口
    \item \textbf{检查代码}：确保代码没有语法错误
    \item \textbf{重启 IDE}：有时候重启可以解决问题
\end{enumerate}

%========== 第二篇：编程基础 ==========
\part{编程基础}

\chapter{Arduino程序结构}

\section{setup()函数}

Arduino 程序有两个核心函数：

\begin{enumerate}
    \item \textbf{setup()}：只运行一次，用于初始化设置
    \item \textbf{loop()}：循环运行，包含主要程序逻辑
\end{enumerate}

setup() 函数通常用于：
\begin{itemize}
    \item 设置引脚模式（输入/输出）
    \item 初始化串口通信
    \item 设置变量初始值
\end{itemize}

\begin{lstlisting}
void setup() {
  // 初始化代码
  pinMode(13, OUTPUT);  // 将引脚13设置为输出模式
  Serial.begin(9600);   // 初始化串口通信
}
\end{lstlisting}

\section{loop()函数}

loop() 函数会不断重复执行，是程序的主要部分：

\begin{lstlisting}
void loop() {
  // 主程序逻辑
  digitalWrite(13, HIGH);  // 点亮LED
  delay(1000);             // 等待1秒
  digitalWrite(13, LOW);   // 熄灭LED
  delay(1000);             // 等待1秒
}
\end{lstlisting}

\section{代码框架}

一个完整的 Arduino 程序结构如下：

\begin{lstlisting}
// 注释：程序说明

// 常量定义
const int ledPin = 13;

// 变量定义
int counter = 0;

void setup() {
  // 初始化设置
  pinMode(ledPin, OUTPUT);
  Serial.begin(9600);
}

void loop() {
  // 主程序逻辑
  digitalWrite(ledPin, HIGH);
  delay(500);
  digitalWrite(ledPin, LOW);
  delay(500);
  
  counter++;
  Serial.println(counter);
}
\end{lstlisting}

\section{注释与代码风格}

注释是代码中非常重要的部分，用于解释代码的作用：

\begin{enumerate}
    \item \textbf{单行注释}：使用 // 开头
    \item \textbf{多行注释}：使用 /* ... */ 包裹
\end{enumerate}

良好的代码风格：
\begin{itemize}
    \item 使用有意义的变量名
    \item 适当添加注释
    \item 保持代码缩进
    \item 一行只写一条语句
\end{itemize}

\chapter{基本语法}

\section{变量与数据类型}

变量是存储数据的容器，Arduino 支持多种数据类型：

\begin{enumerate}
    \item \textbf{int}：整数，范围 -32768 到 32767
    \item \textbf{float}：浮点数，用于存储小数
    \item \textbf{char}：字符，单个字母或符号
    \item \textbf{boolean}：布尔值，true 或 false
    \item \textbf{byte}：无符号字节，范围 0 到 255
\end{enumerate}

变量声明示例：
\begin{lstlisting}
int ledPin = 13;        // 整数变量
float temperature = 25.5; // 浮点数变量
boolean isOn = true;     // 布尔变量
char letter = 'A';       // 字符变量
byte count = 0;          // 字节变量
\end{lstlisting}

\section{常量}

常量是固定不变的值，使用 const 关键字定义：

\begin{lstlisting}
const int LED_PIN = 13;   // 常量，约定用大写字母命名
const float PI = 3.14159; // 圆周率常量
\end{lstlisting}

使用常量的好处：
\begin{itemize}
    \item 提高代码可读性
    \item 便于维护（只需修改一处）
    \item 编译器可以进行优化
\end{itemize}

\section{运算符}

运算符用于执行各种操作：

\begin{enumerate}
    \item \textbf{算术运算符}：+（加）、-（减）、*（乘）、/（除）、\%（取模）
    \item \textbf{比较运算符}：==（等于）、!=（不等于）、>（大于）、<（小于）、>=（大于等于）、<=（小于等于）
    \item \textbf{逻辑运算符}：&&（逻辑与）、||（逻辑或）、!（逻辑非）
    \item \textbf{赋值运算符}：=、+=、-=、*=、/=
\end{enumerate}

示例：
\begin{lstlisting}
int a = 10;
int b = 3;
int c = a + b;      // c = 13
int d = a % b;      // d = 1 (取余数)
boolean result = (a > 5) && (b < 5);  // result = true
\end{lstlisting}

\section{控制语句}

控制语句用于控制程序的执行流程：

\begin{enumerate}
    \item \textbf{if-else 语句}：
    \begin{lstlisting}
if (temperature > 30) {
  digitalWrite(ledPin, HIGH);
} else {
  digitalWrite(ledPin, LOW);
}
    \end{lstlisting}
    
    \item \textbf{for 循环}：
    \begin{lstlisting}
for (int i = 0; i < 10; i++) {
  digitalWrite(13, HIGH);
  delay(100);
  digitalWrite(13, LOW);
  delay(100);
}
    \end{lstlisting}
    
    \item \textbf{while 循环}：
    \begin{lstlisting}
int count = 0;
while (count < 5) {
  Serial.println(count);
  count++;
  delay(1000);
}
    \end{lstlisting}
    
    \item \textbf{switch-case 语句}：
    \begin{lstlisting}
int option = 2;
switch (option) {
  case 1:
    digitalWrite(LED1, HIGH);
    break;
  case 2:
    digitalWrite(LED2, HIGH);
    break;
  default:
    digitalWrite(LED3, HIGH);
}
    \end{lstlisting}
\end{enumerate}

\section{函数}

函数是一段可重复使用的代码块：

\begin{enumerate}
    \item \textbf{定义函数}：
    \begin{lstlisting}
void blinkLED(int pin, int delayTime) {
  digitalWrite(pin, HIGH);
  delay(delayTime);
  digitalWrite(pin, LOW);
  delay(delayTime);
}
    \end{lstlisting}
    
    \item \textbf{调用函数}：
    \begin{lstlisting}
void loop() {
  blinkLED(13, 500);  // 调用函数
}
    \end{lstlisting}
    
    \item \textbf{带返回值的函数}：
    \begin{lstlisting}
int add(int a, int b) {
  return a + b;
}

void loop() {
  int result = add(5, 3);  // result = 8
}
    \end{lstlisting}
\end{enumerate}

\chapter{数字输入输出}

\section{数字信号基础}

数字信号只有两种状态：高电平（HIGH）和低电平（LOW），对应二进制的 1 和 0。

\begin{itemize}
    \item \textbf{HIGH}：高电平，通常是 5V（Arduino UNO）
    \item \textbf{LOW}：低电平，通常是 0V（接地）
\end{itemize}

在 Arduino 中，引脚可以设置为三种模式：
\begin{enumerate}
    \item \textbf{INPUT}：输入模式，用于读取外部信号
    \item \textbf{OUTPUT}：输出模式，用于控制外部设备
    \item \textbf{INPUT_PULLUP}：输入模式（带内部上拉电阻）
\end{enumerate}

\section{pinMode()函数}

用于设置引脚的输入/输出模式：

\begin{lstlisting}
pinMode(pin, mode);
\end{lstlisting}

参数说明：
\begin{itemize}
    \item pin：引脚编号（如 13）
    \item mode：模式（INPUT、OUTPUT 或 INPUT_PULLUP）
\end{itemize}

示例：
\begin{lstlisting}
void setup() {
  pinMode(13, OUTPUT);   // 引脚13设为输出模式
  pinMode(2, INPUT);     // 引脚2设为输入模式
}
\end{lstlisting}

\section{digitalWrite()函数}

用于设置引脚的输出电平：

\begin{lstlisting}
digitalWrite(pin, value);
\end{lstlisting}

参数说明：
\begin{itemize}
    \item pin：引脚编号
    \item value：HIGH（高电平）或 LOW（低电平）
\end{itemize}

示例：
\begin{lstlisting}
void loop() {
  digitalWrite(13, HIGH);  // 点亮LED
  delay(1000);
  digitalWrite(13, LOW);   // 熄灭LED
  delay(1000);
}
\end{lstlisting}

\section{digitalRead()函数}

用于读取引脚的输入电平：

\begin{lstlisting}
int value = digitalRead(pin);
\end{lstlisting}

返回值：
\begin{itemize}
    \item HIGH：高电平
    \item LOW：低电平
\end{itemize}

示例：
\begin{lstlisting}
void setup() {
  pinMode(2, INPUT);    // 引脚2设为输入
  pinMode(13, OUTPUT);  // 引脚13设为输出
}

void loop() {
  int buttonState = digitalRead(2);  // 读取引脚2的状态
  digitalWrite(13, buttonState);    // 根据按钮状态控制LED
}
\end{lstlisting}

\section{实例：LED闪烁}

让 LED 以不同频率闪烁：

\begin{lstlisting}
// 定义LED引脚
const int ledPin = 13;

void setup() {
  // 设置引脚为输出模式
  pinMode(ledPin, OUTPUT);
}

void loop() {
  // 快速闪烁3次
  for (int i = 0; i < 3; i++) {
    digitalWrite(ledPin, HIGH);
    delay(200);
    digitalWrite(ledPin, LOW);
    delay(200);
  }
  
  // 停顿1秒
  delay(1000);
  
  // 慢速闪烁2次
  for (int i = 0; i < 2; i++) {
    digitalWrite(ledPin, HIGH);
    delay(500);
    digitalWrite(ledPin, LOW);
    delay(500);
  }
  
  // 停顿1秒
  delay(1000);
}
\end{lstlisting}

\section{实例：按键控制LED}

使用按键控制 LED 的亮灭：

\begin{lstlisting}
// 定义引脚
const int buttonPin = 2;  // 按键连接到引脚2
const int ledPin = 13;    // LED连接到引脚13

// 变量
int buttonState = 0;      // 存储按键状态

void setup() {
  pinMode(ledPin, OUTPUT);    // LED引脚设为输出
  pinMode(buttonPin, INPUT);  // 按键引脚设为输入
}

void loop() {
  // 读取按键状态
  buttonState = digitalRead(buttonPin);
  
  // 如果按键被按下（高电平）
  if (buttonState == HIGH) {
    digitalWrite(ledPin, HIGH);  // 点亮LED
  } else {
    digitalWrite(ledPin, LOW);   // 熄灭LED
  }
}
\end{lstlisting}

\chapter{模拟输入输出}

\section{模拟信号基础}

模拟信号是连续变化的信号，与数字信号的离散特性不同。

在 Arduino 中：
\begin{itemize}
    \item \textbf{模拟输入}：读取 0-5V 的电压值，转换为 0-1023 的数字
    \item \textbf{模拟输出}：通过 PWM（脉冲宽度调制）实现类似模拟的输出
\end{itemize}

\section{analogRead()函数}

用于读取模拟引脚的电压值：

\begin{lstlisting}
int value = analogRead(pin);
\end{lstlisting}

参数说明：
\begin{itemize}
    \item pin：模拟引脚编号（A0-A5）
\end{itemize}

返回值：0-1023（对应 0-5V）

示例：
\begin{lstlisting}
void setup() {
  Serial.begin(9600);  // 初始化串口
}

void loop() {
  int sensorValue = analogRead(A0);  // 读取A0引脚的模拟值
  Serial.println(sensorValue);       // 打印到串口监视器
  delay(100);
}
\end{lstlisting}

\section{analogWrite()与PWM}

用于输出 PWM 信号：

\begin{lstlisting}
analogWrite(pin, value);
\end{lstlisting}

参数说明：
\begin{itemize}
    \item pin：支持 PWM 的数字引脚（带波浪线标记）
    \item value：0-255（对应占空比 0%-100%）
\end{itemize}

PWM（脉冲宽度调制）通过快速开关来模拟模拟信号。例如，50% 的占空比意味着引脚在一半时间内是高电平，一半时间内是低电平。

示例：
\begin{lstlisting}
void setup() {
  pinMode(9, OUTPUT);  // 引脚9支持PWM
}

void loop() {
  // 从暗到亮
  for (int i = 0; i <= 255; i++) {
    analogWrite(9, i);
    delay(10);
  }
  
  // 从亮到暗
  for (int i = 255; i >= 0; i--) {
    analogWrite(9, i);
    delay(10);
  }
}
\end{lstlisting}

\section{实例：呼吸灯}

让 LED 像呼吸一样渐亮渐暗：

\begin{lstlisting}
// 定义LED引脚（支持PWM）
const int ledPin = 9;

void setup() {
  pinMode(ledPin, OUTPUT);
}

void loop() {
  // 渐亮
  for (int brightness = 0; brightness <= 255; brightness++) {
    analogWrite(ledPin, brightness);
    delay(10);  // 控制变化速度
  }
  
  // 渐暗
  for (int brightness = 255; brightness >= 0; brightness--) {
    analogWrite(ledPin, brightness);
    delay(10);
  }
}
\end{lstlisting}

\section{实例：电位器控制LED亮度}

使用电位器（可变电阻）控制 LED 亮度：

\begin{lstlisting}
// 定义引脚
const int potPin = A0;   // 电位器连接到A0
const int ledPin = 9;    // LED连接到引脚9

void setup() {
  pinMode(ledPin, OUTPUT);
}

void loop() {
  // 读取电位器值（0-1023）
  int potValue = analogRead(potPin);
  
  // 将0-1023映射到0-255
  int brightness = map(potValue, 0, 1023, 0, 255);
  
  // 设置LED亮度
  analogWrite(ledPin, brightness);
  
  delay(50);
}
\end{lstlisting}

map() 函数用于将一个范围的值映射到另一个范围：
\begin{lstlisting}
map(value, fromLow, fromHigh, toLow, toHigh)
\end{lstlisting}

%========== 第三篇：常用模块与传感器 ==========
\part{常用模块与传感器}

\chapter{串口通信}

\section{Serial库介绍}

Serial 库用于 Arduino 与电脑之间的通信，通过 USB 线传输数据。

主要用途：
\begin{itemize}
    \item 调试程序，查看变量值
    \item 发送数据到电脑显示
    \item 接收电脑发送的指令
\end{itemize}

\section{串口输出}

使用 Serial.println() 或 Serial.print() 输出数据：

\begin{lstlisting}
void setup() {
  Serial.begin(9600);  // 初始化串口，波特率为9600
}

void loop() {
  int sensorValue = analogRead(A0);
  
  // 输出数据到串口
  Serial.print("传感器值: ");
  Serial.println(sensorValue);
  
  delay(1000);
}
\end{lstlisting}

打开串口监视器：点击 IDE 右上角的放大镜图标，设置波特率为 9600。

\section{串口读取}

使用 Serial.available() 和 Serial.read() 读取数据：

\begin{lstlisting}
void setup() {
  Serial.begin(9600);
  pinMode(13, OUTPUT);
}

void loop() {
  // 如果有数据可读
  if (Serial.available() > 0) {
    // 读取一个字节
    char incomingByte = Serial.read();
    
    // 根据接收到的字符执行不同操作
    if (incomingByte == 'H') {
      digitalWrite(13, HIGH);
      Serial.println("LED已点亮");
    } else if (incomingByte == 'L') {
      digitalWrite(13, LOW);
      Serial.println("LED已熄灭");
    }
  }
}
\end{lstlisting}

\section{串口事件}

使用 SerialEvent() 处理串口数据：

\begin{lstlisting}
String inputString = "";  // 存储输入的字符串
boolean stringComplete = false;  // 字符串是否接收完成

void setup() {
  Serial.begin(9600);
  inputString.reserve(200);  // 预留空间
}

void loop() {
  if (stringComplete) {
    Serial.println("收到: " + inputString);
    inputString = "";
    stringComplete = false;
  }
}

void serialEvent() {
  while (Serial.available()) {
    char inChar = (char)Serial.read();
    inputString += inChar;
    if (inChar == '\n') {
      stringComplete = true;
    }
  }
}
\end{lstlisting}

\section{实例：串口控制LED}

通过串口发送指令控制 LED：

\begin{lstlisting}
const int ledPin = 13;

void setup() {
  pinMode(ledPin, OUTPUT);
  Serial.begin(9600);
  Serial.println("LED控制程序已启动");
  Serial.println("发送 'on' 点亮LED，发送 'off' 熄灭LED");
}

void loop() {
  if (Serial.available() > 0) {
    String command = Serial.readStringUntil('\n');
    command.trim();  // 去除空格和换行
    
    if (command == "on") {
      digitalWrite(ledPin, HIGH);
      Serial.println("LED已点亮");
    } else if (command == "off") {
      digitalWrite(ledPin, LOW);
      Serial.println("LED已熄灭");
    } else {
      Serial.println("未知命令，请发送 'on' 或 'off'");
    }
  }
}
\end{lstlisting}

\chapter{时间控制}

\section{delay()函数}

delay() 函数用于暂停程序执行指定的毫秒数：

\begin{lstlisting}
delay(ms);
\end{lstlisting}

示例：
\begin{lstlisting}
void loop() {
  digitalWrite(13, HIGH);
  delay(500);  // 暂停500毫秒
  digitalWrite(13, LOW);
  delay(500);
}
\end{lstlisting}

注意：delay() 会阻塞程序，在延迟期间无法执行其他操作。

\section{millis()函数}

millis() 函数返回 Arduino 启动以来的毫秒数，不阻塞程序：

\begin{lstlisting}
unsigned long currentMillis = millis();
\end{lstlisting}

示例：非阻塞延时实现多任务
\begin{lstlisting}
const int led1Pin = 13;
const int led2Pin = 12;

unsigned long previousMillis1 = 0;
unsigned long previousMillis2 = 0;
const long interval1 = 500;  // LED1闪烁间隔
const long interval2 = 1000; // LED2闪烁间隔

void setup() {
  pinMode(led1Pin, OUTPUT);
  pinMode(led2Pin, OUTPUT);
}

void loop() {
  unsigned long currentMillis = millis();
  
  // LED1控制
  if (currentMillis - previousMillis1 >= interval1) {
    previousMillis1 = currentMillis;
    digitalWrite(led1Pin, !digitalRead(led1Pin));
  }
  
  // LED2控制
  if (currentMillis - previousMillis2 >= interval2) {
    previousMillis2 = currentMillis;
    digitalWrite(led2Pin, !digitalRead(led2Pin));
  }
}
\end{lstlisting}

\section{micros()函数}

micros() 函数返回 Arduino 启动以来的微秒数，用于高精度计时：

\begin{lstlisting}
unsigned long currentMicros = micros();
\end{lstlisting}

示例：测量脉冲宽度
\begin{lstlisting}
unsigned long pulseIn(int pin, int value, unsigned long timeout)
\end{lstlisting}

\section{实例：多任务闪烁}

同时控制多个 LED 以不同频率闪烁：

\begin{lstlisting}
const int ledPins[] = {13, 12, 11};  // LED引脚数组
const int ledCount = 3;               // LED数量

unsigned long previousMillis[3] = {0, 0, 0};
const long intervals[3] = {200, 500, 1000};  // 各自的闪烁间隔

void setup() {
  for (int i = 0; i < ledCount; i++) {
    pinMode(ledPins[i], OUTPUT);
  }
}

void loop() {
  unsigned long currentMillis = millis();
  
  for (int i = 0; i < ledCount; i++) {
    if (currentMillis - previousMillis[i] >= intervals[i]) {
      previousMillis[i] = currentMillis;
      digitalWrite(ledPins[i], !digitalRead(ledPins[i]));
    }
  }
}
\end{lstlisting}

\chapter{常用传感器}

\section{温度传感器LM35}

LM35 是一种常用的模拟温度传感器：

特性：
\begin{itemize}
    \item 输出电压与温度成正比
    \item 灵敏度：10mV/°C
    \item 测量范围：-55°C 到 150°C
\end{itemize}

接线方式：
\begin{itemize}
    \item VCC：接 5V
    \item GND：接地
    \item OUT：接模拟引脚（如 A0）
\end{itemize}

示例代码：
\begin{lstlisting}
const int lm35Pin = A0;

void setup() {
  Serial.begin(9600);
}

void loop() {
  int sensorValue = analogRead(lm35Pin);
  
  // 计算电压值（单位：mV）
  float voltage = sensorValue * (5000.0 / 1023.0);
  
  // 计算温度（LM35灵敏度为10mV/°C）
  float temperature = voltage / 10.0;
  
  Serial.print("温度: ");
  Serial.print(temperature);
  Serial.println(" °C");
  
  delay(1000);
}
\end{lstlisting}

\section{光敏电阻}

光敏电阻（LDR）的电阻值随光照强度变化：

接线方式：
\begin{itemize}
    \item 一端接 5V
    \item 另一端接模拟引脚（如 A0）和一个 10kΩ 电阻到 GND
\end{itemize}

示例代码：
\begin{lstlisting}
const int ldrPin = A0;

void setup() {
  Serial.begin(9600);
}

void loop() {
  int ldrValue = analogRead(ldrPin);
  
  // 计算光照强度百分比（假设黑暗时约1023，明亮时约0）
  int brightness = map(ldrValue, 0, 1023, 100, 0);
  
  Serial.print("光照强度: ");
  Serial.print(brightness);
  Serial.println("%");
  
  delay(500);
}
\end{lstlisting}

\section{超声波传感器}

HC-SR04 超声波传感器用于测量距离：

特性：
\begin{itemize}
    \item 测量范围：2cm - 400cm
    \item 精度：±3mm
\end{itemize}

接线方式：
\begin{itemize}
    \item VCC：接 5V
    \item GND：接地
    \item Trig：接数字引脚（如 9）
    \item Echo：接数字引脚（如 10）
\end{itemize}

示例代码：
\begin{lstlisting}
const int trigPin = 9;
const int echoPin = 10;

void setup() {
  Serial.begin(9600);
  pinMode(trigPin, OUTPUT);
  pinMode(echoPin, INPUT);
}

void loop() {
  // 发送触发信号
  digitalWrite(trigPin, LOW);
  delayMicroseconds(2);
  digitalWrite(trigPin, HIGH);
  delayMicroseconds(10);
  digitalWrite(trigPin, LOW);
  
  // 读取回声时间
  long duration = pulseIn(echoPin, HIGH);
  
  // 计算距离（声速约343m/s）
  float distance = duration * 0.0343 / 2;
  
  Serial.print("距离: ");
  Serial.print(distance);
  Serial.println(" cm");
  
  delay(500);
}
\end{lstlisting}

\section{人体红外传感器}

HC-SR501 人体红外传感器用于检测人体移动：

特性：
\begin{itemize}
    \item 检测距离：3-7米可调
    \item 检测角度：110°
\end{itemize}

接线方式：
\begin{itemize}
    \item VCC：接 5V
    \item GND：接地
    \item OUT：接数字引脚（如 2）
\end{itemize}

示例代码：
\begin{lstlisting}
const int sensorPin = 2;
const int ledPin = 13;

void setup() {
  pinMode(sensorPin, INPUT);
  pinMode(ledPin, OUTPUT);
  Serial.begin(9600);
}

void loop() {
  int motion = digitalRead(sensorPin);
  
  if (motion == HIGH) {
    digitalWrite(ledPin, HIGH);
    Serial.println("检测到人体移动");
  } else {
    digitalWrite(ledPin, LOW);
  }
  
  delay(100);
}
\end{lstlisting}

\section{实例：环境监测}

综合使用多种传感器监测环境：

\begin{lstlisting}
// 传感器引脚定义
const int lm35Pin = A0;    // 温度传感器
const int ldrPin = A1;     // 光敏电阻
const int motionPin = 2;   // 人体红外传感器

void setup() {
  pinMode(motionPin, INPUT);
  Serial.begin(9600);
  Serial.println("环境监测系统启动");
}

void loop() {
  // 读取温度
  int tempValue = analogRead(lm35Pin);
  float temperature = (tempValue * 5000.0 / 1023.0) / 10.0;
  
  // 读取光照
  int ldrValue = analogRead(ldrPin);
  int brightness = map(ldrValue, 0, 1023, 100, 0);
  
  // 读取人体检测
  int motion = digitalRead(motionPin);
  
  // 输出数据
  Serial.print("温度: ");
  Serial.print(temperature);
  Serial.print(" °C\t");
  
  Serial.print("光照: ");
  Serial.print(brightness);
  Serial.print("%\t");
  
  Serial.print("人体检测: ");
  Serial.println(motion ? "检测到" : "未检测到");
  
  delay(1000);
}
\end{lstlisting}

\chapter{常用模块}

\section{蜂鸣器模块}

蜂鸣器用于发出声音提示：

接线方式：
\begin{itemize}
    \item VCC：接 5V
    \item GND：接地
    \item I/O：接数字引脚（如 8）
\end{itemize}

示例代码：
\begin{lstlisting}
const int buzzerPin = 8;

void setup() {
  pinMode(buzzerPin, OUTPUT);
}

void loop() {
  // 发出声音
  digitalWrite(buzzerPin, HIGH);
  delay(500);
  
  // 停止发声
  digitalWrite(buzzerPin, LOW);
  delay(500);
}
\end{lstlisting}

播放简单旋律：
\begin{lstlisting}
const int buzzerPin = 8;

// 音符频率定义
#define NOTE_C4  262
#define NOTE_D4  294
#define NOTE_E4  330
#define NOTE_F4  349
#define NOTE_G4  392

void setup() {
  pinMode(buzzerPin, OUTPUT);
}

void playTone(int frequency, int duration) {
  tone(buzzerPin, frequency, duration);
  delay(duration + 50); // 等待播放完成
}

void loop() {
  // 播放简单旋律
  playTone(NOTE_C4, 200);
  playTone(NOTE_D4, 200);
  playTone(NOTE_E4, 200);
  playTone(NOTE_F4, 200);
  playTone(NOTE_G4, 400);
  
  delay(1000);
}
\end{lstlisting}

\section{继电器模块}

继电器用于控制高电压设备：

特性：
\begin{itemize}
    \item 控制端：5V 直流
    \item 被控端：可控制交流 220V 或直流设备
\end{itemize}

接线方式：
\begin{itemize}
    \item VCC：接 5V
    \item GND：接地
    \item IN：接数字引脚（如 7）
\end{itemize}

示例代码：
\begin{lstlisting}
const int relayPin = 7;

void setup() {
  pinMode(relayPin, OUTPUT);
  digitalWrite(relayPin, LOW); // 默认关闭继电器
}

void loop() {
  digitalWrite(relayPin, HIGH); // 打开继电器
  delay(5000);
  
  digitalWrite(relayPin, LOW);  // 关闭继电器
  delay(5000);
}
\end{lstlisting}

\section{舵机控制}

舵机是一种角度可控的电机：

特性：
\begin{itemize}
    \item 角度范围：通常 0°-180°
    \item 控制方式：PWM 信号
\end{itemize}

接线方式：
\begin{itemize}
    \item VCC：接 5V
    \item GND：接地
    \item SIG：接数字引脚（如 9）
\end{itemize}

示例代码：
\begin{lstlisting}
#include <Servo.h>

Servo myservo;  // 创建舵机对象

void setup() {
  myservo.attach(9);  // 将舵机连接到引脚9
}

void loop() {
  // 从0°转到180°
  for (int pos = 0; pos <= 180; pos++) {
    myservo.write(pos);
    delay(15);
  }
  
  // 从180°转到0°
  for (int pos = 180; pos >= 0; pos--) {
    myservo.write(pos);
    delay(15);
  }
}
\end{lstlisting}

\section{红外接收}

红外接收模块用于接收遥控器信号：

接线方式：
\begin{itemize}
    \item VCC：接 5V
    \item GND：接地
    \item OUT：接数字引脚（如 2）
\end{itemize}

示例代码：
\begin{lstlisting}
#include <IRremote.h>

const int RECV_PIN = 2;
IRrecv irrecv(RECV_PIN);
decode_results results;

void setup() {
  Serial.begin(9600);
  irrecv.enableIRIn();  // 启动红外接收
}

void loop() {
  if (irrecv.decode(&results)) {
    Serial.println(results.value, HEX);
    irrecv.resume();  // 接收下一个信号
  }
}
\end{lstlisting}

\section{实例：智能灯控}

使用继电器控制灯光，并通过串口接收指令：

\begin{lstlisting}
const int relayPin = 7;

void setup() {
  pinMode(relayPin, OUTPUT);
  digitalWrite(relayPin, LOW); // 默认关闭
  Serial.begin(9600);
  Serial.println("智能灯控系统启动");
  Serial.println("发送 'on' 开灯，发送 'off' 关灯");
}

void loop() {
  if (Serial.available() > 0) {
    String command = Serial.readStringUntil('\n');
    command.trim();
    
    if (command == "on") {
      digitalWrite(relayPin, HIGH);
      Serial.println("灯已打开");
    } else if (command == "off") {
      digitalWrite(relayPin, LOW);
      Serial.println("灯已关闭");
    } else {
      Serial.println("未知命令");
    }
  }
}
\end{lstlisting}

%========== 第四篇：进阶应用 ==========
\part{进阶应用}

\chapter{中断系统}

\section{什么是中断}

中断是一种硬件机制，可以在特定事件发生时暂停当前程序的执行，转而执行中断服务程序（ISR）。

中断的优点：
\begin{itemize}
    \item 实时响应外部事件
    \item 不阻塞主程序运行
    \item 提高系统效率
\end{itemize}

Arduino UNO 支持外部中断的引脚：
\begin{itemize}
    \item 引脚 2（中断 0）
    \item 引脚 3（中断 1）
\end{itemize}

\section{attachInterrupt()函数}

用于注册中断服务程序：

\begin{lstlisting}
attachInterrupt(digitalPinToInterrupt(pin), ISR, mode);
\end{lstlisting}

参数说明：
\begin{itemize}
    \item pin：引脚编号
    \item ISR：中断服务程序名称
    \item mode：触发模式
\end{itemize}

触发模式：
\begin{enumerate}
    \item \textbf{RISING}：上升沿触发（从 LOW 变为 HIGH）
    \item \textbf{FALLING}：下降沿触发（从 HIGH 变为 LOW）
    \item \textbf{CHANGE}：电平变化时触发
    \item \textbf{LOW}：低电平时触发（不推荐）
\end{enumerate}

示例代码：
\begin{lstlisting}
const int buttonPin = 2;
volatile int buttonPressCount = 0;

void setup() {
  pinMode(buttonPin, INPUT);
  attachInterrupt(digitalPinToInterrupt(buttonPin), buttonISR, RISING);
  Serial.begin(9600);
}

void loop() {
  Serial.print("按键次数: ");
  Serial.println(buttonPressCount);
  delay(500);
}

// 中断服务程序
void buttonISR() {
  buttonPressCount++;
}
\end{lstlisting}

\section{中断使用注意事项}

使用中断时需要注意以下几点：
\begin{itemize}
    \item 中断服务程序应尽量简短
    \item 避免使用 delay() 函数
    \item 访问共享变量时应使用 volatile 关键字
    \item 不要在 ISR 中使用 Serial 通信
    \item 避免在 ISR 中进行复杂运算
\end{itemize}

\section{实例：中断计数}

使用中断实现按键计数功能：

\begin{lstlisting}
const int buttonPin = 2;
const int ledPin = 13;

volatile int count = 0;
volatile boolean ledState = false;

void setup() {
  pinMode(buttonPin, INPUT_PULLUP);  // 使用内部上拉电阻
  pinMode(ledPin, OUTPUT);
  
  // 下降沿触发（按键按下时）
  attachInterrupt(digitalPinToInterrupt(buttonPin), toggleLED, FALLING);
  
  Serial.begin(9600);
}

void loop() {
  Serial.print("计数: ");
  Serial.println(count);
  delay(100);
}

void toggleLED() {
  ledState = !ledState;
  digitalWrite(ledPin, ledState);
  count++;
}
\end{lstlisting}

\chapter{I2C通信}

\section{I2C协议基础}

I2C（Inter-Integrated Circuit）是一种串行通信协议，用于连接多个设备。

特点：
\begin{itemize}
    \item 只需两根线：SDA（数据线）和 SCL（时钟线）
    \item 支持多个设备共享总线
    \item 每个设备有唯一的地址
\end{itemize}

Arduino UNO 的 I2C 引脚：
\begin{itemize}
    \item SDA：A4
    \item SCL：A5
\end{itemize}

\section{Wire库介绍}

Wire 库用于实现 I2C 通信：

主要函数：
\begin{enumerate}
    \item \textbf{Wire.begin()}：初始化 I2C 总线
    \item \textbf{Wire.beginTransmission(address)}：开始向指定地址发送数据
    \item \textbf{Wire.write(data)}：写入数据
    \item \textbf{Wire.endTransmission()}：结束传输
    \item \textbf{Wire.requestFrom(address, quantity)}：从指定地址请求数据
    \item \textbf{Wire.available()}：检查是否有数据可读
    \item \textbf{Wire.read()}：读取一个字节
\end{enumerate}

\section{I2C设备地址扫描}

扫描总线上连接的 I2C 设备：

\begin{lstlisting}
#include <Wire.h>

void setup() {
  Wire.begin();
  Serial.begin(9600);
  Serial.println("I2C 设备扫描...");
  
  byte error, address;
  int nDevices = 0;
  
  for (address = 1; address < 127; address++) {
    Wire.beginTransmission(address);
    error = Wire.endTransmission();
    
    if (error == 0) {
      Serial.print("找到设备地址: 0x");
      if (address < 16) Serial.print("0");
      Serial.println(address, HEX);
      nDevices++;
    }
  }
  
  if (nDevices == 0) {
    Serial.println("未找到任何I2C设备");
  }
}

void loop() {}
\end{lstlisting}

\section{实例：I2C LCD显示屏}

使用 I2C 接口的 16x2 LCD 显示屏：

\begin{lstlisting}
#include <Wire.h>
#include <LiquidCrystal_I2C.h>

// 设置 LCD 地址和尺寸
LiquidCrystal_I2C lcd(0x27, 16, 2);

void setup() {
  // 初始化 LCD
  lcd.init();
  // 开启背光
  lcd.backlight();
  
  // 显示欢迎信息
  lcd.setCursor(0, 0);
  lcd.print("Hello, Arduino!");
  lcd.setCursor(0, 1);
  lcd.print("I2C LCD Test");
}

void loop() {
  // 显示当前时间
  lcd.setCursor(0, 1);
  lcd.print(millis() / 1000);
  lcd.print("s");
  delay(1000);
}
\end{lstlisting}

\chapter{SPI通信}

\section{SPI协议基础}

SPI（Serial Peripheral Interface）是一种高速串行通信协议。

特点：
\begin{itemize}
    \item 全双工通信
    \item 速度比 I2C 快
    \item 需要更多引脚
\end{itemize}

SPI 总线包含四根线：
\begin{itemize}
    \item SCK：时钟线
    \item MOSI：主设备输出，从设备输入
    \item MISO：主设备输入，从设备输出
    \item SS/CS：从设备选择线
\end{itemize}

Arduino UNO 的 SPI 引脚：
\begin{itemize}
    \item SCK：13
    \item MOSI：11
    \item MISO：12
    \item SS：10（默认，但可使用任意数字引脚）
\end{itemize}

\section{SPI库介绍}

SPI 库用于实现 SPI 通信：

主要函数：
\begin{enumerate}
    \item \textbf{SPI.begin()}：初始化 SPI 总线
    \item \textbf{SPI.beginTransaction(SPISettings(maxSpeed, dataOrder, dataMode))}：开始 SPI 事务
    \item \textbf{SPI.transfer(data)}：发送并接收一个字节
    \item \textbf{SPI.endTransaction()}：结束 SPI 事务
    \item \textbf{SPI.end()}：禁用 SPI 总线
\end{enumerate}

SPISettings 参数：
\begin{itemize}
    \item maxSpeed：最大时钟频率
    \item dataOrder：MSBFIRST（高位在前）或 LSBFIRST（低位在前）
    \item dataMode：SPI 模式（0-3），决定时钟极性和相位
\end{itemize}

\section{实例：SPI通信}

SPI 主设备发送数据到从设备：

主设备代码：
\begin{lstlisting}
#include <SPI.h>

const int slaveSelectPin = 10;

void setup() {
  pinMode(slaveSelectPin, OUTPUT);
  digitalWrite(slaveSelectPin, HIGH); // 禁用从设备
  SPI.begin();
  Serial.begin(9600);
}

void loop() {
  digitalWrite(slaveSelectPin, LOW); // 启用从设备
  
  // 发送数据
  byte data = 0x55;
  byte received = SPI.transfer(data);
  
  digitalWrite(slaveSelectPin, HIGH); // 禁用从设备
  
  Serial.print("发送: ");
  Serial.println(data, HEX);
  Serial.print("接收: ");
  Serial.println(received, HEX);
  
  delay(1000);
}
\end{lstlisting}

从设备代码：
\begin{lstlisting}
#include <SPI.h>

const int slaveSelectPin = 10;

void setup() {
  pinMode(slaveSelectPin, INPUT);
  pinMode(MISO, OUTPUT);
  SPCR |= _BV(SPE); // 启用 SPI 从设备模式
  Serial.begin(9600);
}

void loop() {
  if (!(SPSR & _BV(SPIF))) return; // 等待数据
  
  byte received = SPDR; // 读取接收到的数据
  byte response = received + 1; // 处理数据
  SPDR = response; // 设置响应数据
  
  Serial.print("从设备接收: ");
  Serial.println(received, HEX);
}
\end{lstlisting}

\chapter{EEPROM使用}

\section{EEPROM库介绍}

EEPROM（Electrically Erasable Programmable Read-Only Memory）是一种非易失性存储器，可以在断电后保存数据。

Arduino UNO 的 EEPROM 容量：
\begin{itemize}
    \item ATmega328P：1KB（1024 字节）
    \item 地址范围：0-1023
\end{itemize}

使用 EEPROM 库需要先包含头文件：
\begin{lstlisting}
#include <EEPROM.h>
\end{lstlisting}

\section{数据读写}

主要函数：
\begin{enumerate}
    \item \textbf{EEPROM.write(address, value)}：向指定地址写入字节
    \item \textbf{EEPROM.read(address)}：从指定地址读取字节
    \item \textbf{EEPROM.update(address, value)}：仅当值不同时才写入
    \item \textbf{EEPROM.get(address, data)}：读取任意类型的数据
    \item \textbf{EEPROM.put(address, data)}：写入任意类型的数据
\end{enumerate}

基本读写示例：
\begin{lstlisting}
#include <EEPROM.h>

void setup() {
  // 写入数据
  EEPROM.write(0, 123);      // 在地址0写入字节
  EEPROM.write(1, 255);      // 在地址1写入字节
  
  // 读取数据
  byte value1 = EEPROM.read(0);
  byte value2 = EEPROM.read(1);
  
  Serial.begin(9600);
  Serial.print("地址0的值: ");
  Serial.println(value1);
  Serial.print("地址1的值: ");
  Serial.println(value2);
}

void loop() {}
\end{lstlisting}

读写任意类型数据：
\begin{lstlisting}
#include <EEPROM.h>

struct Config {
  int brightness;
  float temperature;
  char name[20];
};

void setup() {
  Config config;
  
  // 写入配置
  config.brightness = 128;
  config.temperature = 25.5;
  strcpy(config.name, "Arduino");
  
  EEPROM.put(0, config);
  
  // 读取配置
  Config readConfig;
  EEPROM.get(0, readConfig);
  
  Serial.begin(9600);
  Serial.print("亮度: ");
  Serial.println(readConfig.brightness);
  Serial.print("温度: ");
  Serial.println(readConfig.temperature);
  Serial.print("名称: ");
  Serial.println(readConfig.name);
}

void loop() {}
\end{lstlisting}

\section{实例：保存配置信息}

保存用户配置信息到 EEPROM：

\begin{lstlisting}
#include <EEPROM.h>

// 配置结构体
struct DeviceConfig {
  int ledBrightness;
  int delayTime;
  boolean autoMode;
};

const int CONFIG_VERSION = 1;
const int CONFIG_ADDRESS = 0;

void setup() {
  Serial.begin(9600);
  
  // 检查配置版本
  int storedVersion = EEPROM.read(CONFIG_ADDRESS);
  
  if (storedVersion != CONFIG_VERSION) {
    // 首次运行或版本升级，写入默认配置
    Serial.println("首次运行，写入默认配置");
    DeviceConfig defaultConfig = {128, 1000, true};
    EEPROM.write(CONFIG_ADDRESS, CONFIG_VERSION);
    EEPROM.put(CONFIG_ADDRESS + 1, defaultConfig);
  }
  
  // 读取配置
  DeviceConfig config;
  EEPROM.get(CONFIG_ADDRESS + 1, config);
  
  Serial.println("当前配置:");
  Serial.print("LED亮度: ");
  Serial.println(config.ledBrightness);
  Serial.print("延时时间: ");
  Serial.println(config.delayTime);
  Serial.print("自动模式: ");
  Serial.println(config.autoMode ? "开启" : "关闭");
}

void loop() {}
\end{lstlisting}

注意事项：
\begin{itemize}
    \item EEPROM 有写入次数限制（约 10 万次）
    \item 避免频繁写入
    \item 使用 update() 代替 write() 可延长寿命
\end{itemize}

%========== 第五篇：综合项目 ==========
\part{综合项目}

\chapter{项目一：LED灯控制}

\section{项目需求}

本项目实现一个智能LED灯控制系统，具有以下功能：
\begin{itemize}
    \item 通过按键控制LED的开关
    \item 通过电位器调节LED亮度
    \item 支持自动呼吸灯模式
    \item 通过串口接收指令控制
\end{itemize}

\section{硬件准备}

所需材料：
\begin{enumerate}
    \item Arduino UNO 开发板
    \item LED 灯（红色）
    \item 220Ω 电阻
    \item 10kΩ 电位器
    \item 按键
    \item 面包板和杜邦线
\end{enumerate}

\section{电路搭建}

连接方式：
\begin{enumerate}
    \item LED 正极 → 220Ω 电阻 → 数字引脚 9（PWM）
    \item LED 负极 → GND
    \item 电位器一端 → 5V，一端 → GND，中间 → A0
    \item 按键一端 → 数字引脚 2，一端 → GND
    \item 使用内部上拉电阻
\end{enumerate}

\section{代码实现}

完整代码：
\begin{lstlisting}
// 引脚定义
const int ledPin = 9;
const int potPin = A0;
const int buttonPin = 2;

// 状态变量
boolean ledState = false;
boolean autoMode = false;
unsigned long previousMillis = 0;
int brightness = 0;
int fadeDirection = 1;

void setup() {
  pinMode(ledPin, OUTPUT);
  pinMode(buttonPin, INPUT_PULLUP);
  Serial.begin(9600);
  
  // 注册中断
  attachInterrupt(digitalPinToInterrupt(buttonPin), toggleLED, FALLING);
  
  Serial.println("LED控制系统启动");
  Serial.println("指令：on-开灯 off-关灯 auto-自动模式");
}

void loop() {
  // 处理串口指令
  if (Serial.available() > 0) {
    String command = Serial.readStringUntil('\n');
    command.trim();
    
    if (command == "on") {
      ledState = true;
      autoMode = false;
      analogWrite(ledPin, 255);
      Serial.println("LED已开启");
    } else if (command == "off") {
      ledState = false;
      autoMode = false;
      analogWrite(ledPin, 0);
      Serial.println("LED已关闭");
    } else if (command == "auto") {
      autoMode = !autoMode;
      Serial.println(autoMode ? "自动模式已开启" : "自动模式已关闭");
    }
  }
  
  // 自动呼吸灯模式
  if (autoMode) {
    unsigned long currentMillis = millis();
    if (currentMillis - previousMillis >= 10) {
      previousMillis = currentMillis;
      
      brightness += fadeDirection;
      
      if (brightness >= 255) {
        fadeDirection = -1;
      } else if (brightness <= 0) {
        fadeDirection = 1;
      }
      
      analogWrite(ledPin, brightness);
    }
  } else if (ledState) {
    // 手动模式：电位器控制亮度
    int potValue = analogRead(potPin);
    brightness = map(potValue, 0, 1023, 0, 255);
    analogWrite(ledPin, brightness);
  }
}

// 中断服务程序：按键切换
void toggleLED() {
  if (!autoMode) {
    ledState = !ledState;
    if (ledState) {
      analogWrite(ledPin, brightness > 0 ? brightness : 255);
    } else {
      analogWrite(ledPin, 0);
    }
  }
}
\end{lstlisting}

\section{项目总结}

本项目综合运用了以下知识点：
\begin{enumerate}
    \item 数字输出和 PWM 控制
    \item 模拟输入读取
    \item 外部中断使用
    \item 串口通信
    \item 状态机设计
\end{enumerate}

\chapter{项目二：温湿度监测}

\section{项目需求}

本项目实现一个温湿度监测系统，具有以下功能：
\begin{itemize}
    \item 实时监测温度和湿度
    \item 在 LCD 显示屏上显示数据
    \item 温度或湿度超过阈值时报警
    \item 通过串口上传数据
\end{itemize}

\section{硬件准备}

所需材料：
\begin{enumerate}
    \item Arduino UNO 开发板
    \item DHT11 温湿度传感器
    \item 16x2 LCD 显示屏（I2C接口）
    \item 蜂鸣器模块
    \item 面包板和杜邦线
\end{enumerate}

\section{电路搭建}

连接方式：
\begin{enumerate}
    \item DHT11 VCC → 5V
    \item DHT11 GND → GND
    \item DHT11 DATA → 数字引脚 2
    \item LCD SDA → A4
    \item LCD SCL → A5
    \item 蜂鸣器 VCC → 5V
    \item 蜂鸣器 GND → GND
    \item 蜂鸣器 I/O → 数字引脚 8
\end{enumerate}

\section{代码实现}

完整代码：
\begin{lstlisting}
#include <DHT.h>
#include <LiquidCrystal_I2C.h>

// 传感器设置
#define DHTPIN 2
#define DHTTYPE DHT11
DHT dht(DHTPIN, DHTTYPE);

// LCD 设置
LiquidCrystal_I2C lcd(0x27, 16, 2);

// 蜂鸣器引脚
const int buzzerPin = 8;

// 阈值设置
const float tempThreshold = 30.0;
const float humiThreshold = 70.0;

void setup() {
  pinMode(buzzerPin, OUTPUT);
  digitalWrite(buzzerPin, LOW);
  
  // 初始化传感器
  dht.begin();
  
  // 初始化 LCD
  lcd.init();
  lcd.backlight();
  lcd.print("温湿度监测系统");
  delay(2000);
  
  // 初始化串口
  Serial.begin(9600);
}

void loop() {
  // 读取传感器数据
  float humidity = dht.readHumidity();
  float temperature = dht.readTemperature();
  
  // 检查读取是否成功
  if (isnan(humidity) || isnan(temperature)) {
    lcd.clear();
    lcd.setCursor(0, 0);
    lcd.print("读取失败!");
    Serial.println("传感器读取失败");
    delay(1000);
    return;
  }
  
  // 显示数据到 LCD
  lcd.clear();
  lcd.setCursor(0, 0);
  lcd.print("温度: ");
  lcd.print(temperature);
  lcd.print(" C");
  
  lcd.setCursor(0, 1);
  lcd.print("湿度: ");
  lcd.print(humidity);
  lcd.print(" %");
  
  // 串口输出
  Serial.print("温度: ");
  Serial.print(temperature);
  Serial.print(" C, 湿度: ");
  Serial.print(humidity);
  Serial.println(" %");
  
  // 报警检测
  boolean alarm = false;
  if (temperature > tempThreshold) {
    alarm = true;
    lcd.setCursor(12, 0);
    lcd.print("!");
  }
  if (humidity > humiThreshold) {
    alarm = true;
    lcd.setCursor(12, 1);
    lcd.print("!");
  }
  
  // 控制蜂鸣器
  digitalWrite(buzzerPin, alarm ? HIGH : LOW);
  
  delay(2000);
}
\end{lstlisting}

\section{项目总结}

本项目综合运用了以下知识点：
\begin{enumerate}
    \item 传感器数据读取
    \item I2C 通信
    \item LCD 显示
    \item 报警系统设计
    \item 串口数据上传
\end{enumerate}

\chapter{项目三：智能小车}

\section{项目需求}

本项目实现一个遥控智能小车，具有以下功能：
\begin{itemize}
    \item 通过串口接收控制指令
    \item 支持前进、后退、左转、右转
    \item 使用超声波传感器避障
    \item 支持自动避障模式
\end{itemize}

\section{硬件准备}

所需材料：
\begin{enumerate}
    \item Arduino UNO 开发板
    \item L298N 电机驱动模块
    \item 直流电机（带减速箱）× 2
    \item 小车底盘和车轮
    \item HC-SR04 超声波传感器
    \item 面包板和杜邦线
    \item 9V 电池组
\end{enumerate}

\section{电路搭建}

连接方式：
\begin{enumerate}
    \item L298N VCC → 9V 电池
    \item L298N GND → Arduino GND
    \item L298N 5V → Arduino 5V（可选）
    \item L298N IN1 → 数字引脚 8
    \item L298N IN2 → 数字引脚 9
    \item L298N IN3 → 数字引脚 10
    \item L298N IN4 → 数字引脚 11
    \item HC-SR04 Trig → 数字引脚 12
    \item HC-SR04 Echo → 数字引脚 13
\end{enumerate}

\section{代码实现}

完整代码：
\begin{lstlisting}
// 电机控制引脚
const int IN1 = 8;
const int IN2 = 9;
const int IN3 = 10;
const int IN4 = 11;

// 超声波传感器引脚
const int trigPin = 12;
const int echoPin = 13;

// 自动避障标志
boolean autoMode = false;

void setup() {
  // 设置电机引脚为输出
  pinMode(IN1, OUTPUT);
  pinMode(IN2, OUTPUT);
  pinMode(IN3, OUTPUT);
  pinMode(IN4, OUTPUT);
  
  // 设置传感器引脚
  pinMode(trigPin, OUTPUT);
  pinMode(echoPin, INPUT);
  
  // 初始化串口
  Serial.begin(9600);
  Serial.println("智能小车控制系统");
  Serial.println("指令: f-前进 b-后退 l-左转 r-右转 s-停止 a-自动模式");
  
  // 初始停止
  stop();
}

void loop() {
  // 处理串口指令
  if (Serial.available() > 0) {
    char cmd = Serial.read();
    
    switch(cmd) {
      case 'f':
        autoMode = false;
        forward();
        break;
      case 'b':
        autoMode = false;
        backward();
        break;
      case 'l':
        autoMode = false;
        left();
        break;
      case 'r':
        autoMode = false;
        right();
        break;
      case 's':
        autoMode = false;
        stop();
        break;
      case 'a':
        autoMode = !autoMode;
        Serial.println(autoMode ? "自动避障模式" : "手动控制模式");
        break;
    }
  }
  
  // 自动避障模式
  if (autoMode) {
    float distance = getDistance();
    Serial.print("距离: ");
    Serial.println(distance);
    
    if (distance < 20) {
      // 距离过近，后退并转向
      backward();
      delay(500);
      left();
      delay(500);
    } else {
      forward();
    }
  }
}

// 获取距离
float getDistance() {
  digitalWrite(trigPin, LOW);
  delayMicroseconds(2);
  digitalWrite(trigPin, HIGH);
  delayMicroseconds(10);
  digitalWrite(trigPin, LOW);
  
  long duration = pulseIn(echoPin, HIGH);
  return duration * 0.0343 / 2;
}

// 前进
void forward() {
  digitalWrite(IN1, HIGH);
  digitalWrite(IN2, LOW);
  digitalWrite(IN3, HIGH);
  digitalWrite(IN4, LOW);
  Serial.println("前进");
}

// 后退
void backward() {
  digitalWrite(IN1, LOW);
  digitalWrite(IN2, HIGH);
  digitalWrite(IN3, LOW);
  digitalWrite(IN4, HIGH);
  Serial.println("后退");
}

// 左转
void left() {
  digitalWrite(IN1, LOW);
  digitalWrite(IN2, HIGH);
  digitalWrite(IN3, HIGH);
  digitalWrite(IN4, LOW);
  Serial.println("左转");
}

// 右转
void right() {
  digitalWrite(IN1, HIGH);
  digitalWrite(IN2, LOW);
  digitalWrite(IN3, LOW);
  digitalWrite(IN4, HIGH);
  Serial.println("右转");
}

// 停止
void stop() {
  digitalWrite(IN1, LOW);
  digitalWrite(IN2, LOW);
  digitalWrite(IN3, LOW);
  digitalWrite(IN4, LOW);
  Serial.println("停止");
}
\end{lstlisting}

\section{项目总结}

本项目综合运用了以下知识点：
\begin{enumerate}
    \item 电机驱动控制
    \item 超声波测距
    \item 串口通信
    \item 自动避障算法
    \item 状态机设计
\end{enumerate}

\chapter{项目四：物联网应用}

\section{项目需求}

本项目实现一个基于 WiFi 的物联网温湿度监测系统，具有以下功能：
\begin{itemize}
    \item 通过 WiFi 连接到互联网
    \item 定时上传温湿度数据到云平台
    \item 支持远程控制 LED
    \item 接收云平台指令
\end{itemize}

\section{硬件准备}

所需材料：
\begin{enumerate}
    \item Arduino UNO 开发板
    \item ESP8266 WiFi 模块
    \item DHT11 温湿度传感器
    \item LED 灯和电阻
    \item 面包板和杜邦线
\end{enumerate}

\section{电路搭建}

连接方式：
\begin{enumerate}
    \item ESP8266 VCC → 3.3V
    \item ESP8266 GND → GND
    \item ESP8266 TX → Arduino RX（引脚0）
    \item ESP8266 RX → Arduino TX（引脚1）
    \item DHT11 VCC → 5V
    \item DHT11 GND → GND
    \item DHT11 DATA → 数字引脚 2
    \item LED → 数字引脚 13
\end{enumerate}

\section{代码实现}

完整代码：
\begin{lstlisting}
#include <SoftwareSerial.h>
#include <DHT.h>

// WiFi模块串口
SoftwareSerial espSerial(2, 3); // RX, TX

// DHT传感器
#define DHTPIN 4
#define DHTTYPE DHT11
DHT dht(DHTPIN, DHTTYPE);

// LED引脚
const int ledPin = 13;

// WiFi配置
const char* ssid = "YOUR_WIFI_SSID";
const char* password = "YOUR_WIFI_PASSWORD";
const char* server = "api.example.com";

void setup() {
  pinMode(ledPin, OUTPUT);
  digitalWrite(ledPin, LOW);
  
  // 初始化串口
  Serial.begin(9600);
  espSerial.begin(115200);
  
  // 初始化传感器
  dht.begin();
  
  delay(1000);
  
  // 连接WiFi
  connectWiFi();
}

void loop() {
  // 读取传感器数据
  float humidity = dht.readHumidity();
  float temperature = dht.readTemperature();
  
  // 上传数据
  sendData(temperature, humidity);
  
  // 检查指令
  checkCommand();
  
  delay(5000);
}

// 连接WiFi
void connectWiFi() {
  espSerial.println("AT+CWMODE=1");
  delay(1000);
  
  String cmd = "AT+CWJAP=\"" + String(ssid) + "\",\"" + String(password) + "\"";
  espSerial.println(cmd);
  delay(5000);
  
  espSerial.println("AT+CIPMUX=1");
  delay(1000);
}

// 发送数据
void sendData(float temp, float humi) {
  espSerial.println("AT+CIPSTART=0,\"TCP\",\"" + String(server) + "\",80");
  delay(2000);
  
  String data = "GET /upload?temp=" + String(temp) + "&humi=" + String(humi) + " HTTP/1.1\r\n";
  data += "Host: " + String(server) + "\r\n";
  data += "Connection: close\r\n\r\n";
  
  espSerial.println("AT+CIPSEND=0," + String(data.length()));
  delay(1000);
  
  espSerial.print(data);
  delay(2000);
  
  espSerial.println("AT+CIPCLOSE=0");
  delay(1000);
}

// 检查指令
void checkCommand() {
  if (espSerial.available() > 0) {
    String response = espSerial.readStringUntil('\n');
    
    if (response.indexOf("LED_ON") >= 0) {
      digitalWrite(ledPin, HIGH);
      Serial.println("LED已开启");
    } else if (response.indexOf("LED_OFF") >= 0) {
      digitalWrite(ledPin, LOW);
      Serial.println("LED已关闭");
    }
  }
}
\end{lstlisting}

\section{项目总结}

本项目综合运用了以下知识点：
\begin{enumerate}
    \item WiFi 模块控制
    \item HTTP 协议通信
    \item 传感器数据采集
    \item 云平台对接
    \item 远程控制
\end{enumerate}

%========== 第六篇：附录 ==========
\part{附录}

\chapter{常用函数速查}

\section{数字I/O函数}

\begin{itemize}
    \item \textbf{pinMode(pin, mode)}：设置引脚模式（INPUT/OUTPUT/INPUT\_PULLUP）
    \item \textbf{digitalWrite(pin, value)}：设置引脚输出电平（HIGH/LOW）
    \item \textbf{digitalRead(pin)}：读取引脚输入电平
\end{itemize}

\section{模拟I/O函数}

\begin{itemize}
    \item \textbf{analogRead(pin)}：读取模拟值（0-1023）
    \item \textbf{analogWrite(pin, value)}：输出PWM信号（0-255）
    \item \textbf{map(value, fromLow, fromHigh, toLow, toHigh)}：映射数值范围
\end{itemize}

\section{高级I/O函数}

\begin{itemize}
    \item \textbf{pulseIn(pin, value, timeout)}：测量脉冲宽度
    \item \textbf{shiftOut(dataPin, clockPin, bitOrder, value)}：移位输出
    \item \textbf{shiftIn(dataPin, clockPin, bitOrder)}：移位输入
    \item \textbf{tone(pin, frequency, duration)}：产生音调
    \item \textbf{noTone(pin)}：停止音调
\end{itemize}

\section{时间函数}

\begin{itemize}
    \item \textbf{delay(ms)}：延时指定毫秒数
    \item \textbf{delayMicroseconds(us)}：延时指定微秒数
    \item \textbf{millis()}：返回系统运行毫秒数
    \item \textbf{micros()}：返回系统运行微秒数
\end{itemize}

\section{数学函数}

\begin{itemize}
    \item \textbf{min(a, b)}：返回较小值
    \item \textbf{max(a, b)}：返回较大值
    \item \textbf{abs(x)}：返回绝对值
    \item \textbf{constrain(x, min, max)}：限制数值范围
    \item \textbf{random(min, max)}：生成随机数
    \item \textbf{sqrt(x)}：计算平方根
    \item \textbf{sin(rad)}：正弦函数
    \item \textbf{cos(rad)}：余弦函数
\end{itemize}

\section{串口函数}

\begin{itemize}
    \item \textbf{Serial.begin(baud)}：初始化串口通信
    \item \textbf{Serial.print(data)}：输出数据
    \item \textbf{Serial.println(data)}：输出数据并换行
    \item \textbf{Serial.available()}：返回可读取字节数
    \item \textbf{Serial.read()}：读取一个字节
    \item \textbf{Serial.readString()}：读取字符串
\end{itemize}

\chapter{常见问题与解答}

\section{编译错误}

\begin{enumerate}
    \item \textbf{问题："'Serial' was not declared in this scope"}
    \\解答：需要在 setup() 中调用 Serial.begin()
    
    \item \textbf{问题："expected '}' at end of input"}
    \\解答：检查代码中是否有缺少的大括号
    
    \item \textbf{问题："invalid conversion from 'int' to 'const char*'"}
    \\解答：使用 String() 函数将整数转换为字符串
    
    \item \textbf{问题："'delay' was not declared in this scope"}
    \\解答：delay() 函数应该在 loop() 或其他函数内部调用
\end{enumerate}

\section{上传错误}

\begin{enumerate}
    \item \textbf{问题："port busy" 或 "permission denied"}
    \\解答：关闭其他可能占用串口的程序，如串口监视器、其他IDE等
    
    \item \textbf{问题："avrdude: stk500_recv(): programmer is not responding"}
    \\解答：
    \begin{itemize}
        \item 检查 USB 连接是否正常
        \item 确认选择了正确的开发板和端口
        \item 尝试按下开发板上的复位按钮
    \end{itemize}
    
    \item \textbf{问题："board not found"}
    \\解答：确认安装了正确的开发板支持包
\end{enumerate}

\section{运行时错误}

\begin{enumerate}
    \item \textbf{问题：LED不亮}
    \\解答：
    \begin{itemize}
        \item 检查电路连接是否正确
        \item 确认引脚模式设置正确（OUTPUT）
        \item 使用数字万用表测量引脚电压
    \end{itemize}
    
    \item \textbf{问题：传感器读数异常}
    \\解答：
    \begin{itemize}
        \item 检查传感器接线是否正确
        \item 确认传感器供电电压是否符合要求
        \item 检查传感器模块是否损坏
    \end{itemize}
    
    \item \textbf{问题：程序运行速度异常}
    \\解答：检查是否有过长的 delay() 调用，考虑使用 millis() 替代
\end{enumerate}

\section{硬件故障}

\begin{enumerate}
    \item \textbf{问题：开发板无法通电}
    \\解答：
    \begin{itemize}
        \item 检查 USB 线是否正常
        \item 确认电源适配器输出是否正确
        \item 检查开发板是否短路
    \end{itemize}
    
    \item \textbf{问题：引脚无输出}
    \\解答：
    \begin{itemize}
        \item 使用数字万用表测量引脚电平
        \item 确认引脚是否被正确配置为 OUTPUT
        \item 检查是否有外部电路短路
    \end{itemize}
    
    \item \textbf{问题：串口无数据}
    \\解答：
    \begin{itemize}
        \item 确认串口监视器波特率设置正确
        \item 检查 USB 连接
        \item 确认代码中调用了 Serial.begin()
    \end{itemize}
\end{enumerate}

\chapter{参考资料}

\section{Arduino官方文档}

\begin{itemize}
    \item Arduino 官网：https://www.arduino.cc
    \item Arduino 参考文档：https://www.arduino.cc/reference/zh
    \item Arduino 教程：https://www.arduino.cc/en/Tutorial/HomePage
\end{itemize}

\section{推荐学习资源}

\begin{itemize}
    \item \textbf{书籍}：
    \begin{itemize}
        \item 《Arduino编程从零开始》
        \item 《Arduino实战指南》
        \item 《爱上Arduino》
    \end{itemize}
    
    \item \textbf{在线教程}：
    \begin{itemize}
        \item Arduino 官方教程：https://www.arduino.cc/en/Tutorial
        \item Adafruit 教程：https://learn.adafruit.com
        \item SparkFun 教程：https://learn.sparkfun.com
    \end{itemize}
    
    \item \textbf{社区论坛}：
    \begin{itemize}
        \item Arduino 中文论坛：https://www.arduino.cn
        \item Stack Overflow：https://stackoverflow.com/questions/tagged/arduino
    \end{itemize}
\end{itemize}

\section{常用开发板规格}

\begin{enumerate}
    \item \textbf{Arduino UNO}：
    \begin{itemize}
        \item 控制器：ATmega328P
        \item 数字I/O：14个（6个PWM）
        \item 模拟输入：6个
        \item 内存：32KB Flash，2KB SRAM，1KB EEPROM
    \end{itemize}
    
    \item \textbf{Arduino Nano}：
    \begin{itemize}
        \item 控制器：ATmega328P
        \item 数字I/O：14个（6个PWM）
        \item 模拟输入：8个
        \item 内存：32KB Flash，2KB SRAM，1KB EEPROM
    \end{itemize}
    
    \item \textbf{Arduino Mega}：
    \begin{itemize}
        \item 控制器：ATmega2560
        \item 数字I/O：54个（15个PWM）
        \item 模拟输入：16个
        \item 内存：256KB Flash，8KB SRAM，4KB EEPROM
    \end{itemize}
    
    \item \textbf{Arduino Leonardo}：
    \begin{itemize}
        \item 控制器：ATmega32U4
        \item 数字I/O：20个（7个PWM）
        \item 模拟输入：12个
        \item 内存：32KB Flash，2.5KB SRAM，1KB EEPROM
        \item 支持 USB HID 功能
    \end{itemize}
\end{enumerate}

\backmatter

\end{document}
